\documentclass[12pt, twoside]{article}
% Emit local LDF shims before packages load
\input{lib/datatool-ldf}

% Make it possible to conditionally depend on the TeX engine used
\RequirePackage{ifluatex}

\RequirePackage{comment}

% To access the \captionof{} functionality
\RequirePackage{caption}

\RequirePackage{luacode}

%\RequirePackage{silence}
%\WarningFilter{latexfont}{Font shape `U/stmry/m/n' in size <5.5> not available}
%\WarningFilter{latexfont}{Font shape}

\newif\ifinswedish
\inswedishfalse

\RequirePackage{etoolbox}

% The following toggles are needed because of kth/kthcolors.tex and lib/includes.tex
\newif\ifnomenclature
\nomenclaturefalse

\newif\ifdigitaloutput
\digitaloutputfalse

\makeatletter
%% Define a pair of commands to disable and reenable specific packages - see https://tex.stackexchange.com/questions/39415/unload-a-latex-package
\makeatletter
\newcommand{\disablepackage}[2]{%
  \disable@package@load{#1}{#2}%
}
\newcommand{\reenablepackage}[1]{%
  \reenable@package@load{#1}%
}
\makeatother

%% To avoid the warning: "Package transparent Warning: Loading aborted, because pdfTeX is not running in PDF mode."
\disablepackage{transparent}{}

% include a variety of packages that are useful
\input{lib/includes}

\usepackage{minted}
\RequirePackage{fancyhdr} % Take control of headers and footers

\usepackage{luatexja-fontspec}
%\ltjdefcharrange{8}{"2190-"21FF} % Define range 8 as Arrows
%\ltjsetparameter{jacharrange={-8}} % Treat range 8 as Western characters (ALChar)
\ltjsetparameter{jacharrange={-3}} % Treat range 3 as Western characters (ALChar)

\begin{comment}
\ifinswedish
    \usepackage[english, main=swedish, bidi=basic]{babel}
\else
    \usepackage[swedish, main=english, provide+=*, bidi=basic]{babel}
\fi
\end{comment}
\ifinswedish
    \usepackage[main=swedish]{babel}
\else
    \usepackage[main=english, provide+=*]{babel}
\fi
\setmainjfont{Noto Sans CJK JP}
% Declare STIX Two Math or Noto Sans Symbols as a fallback for Japanese symbols
%\ltjdeclarealtfontrange{arrows}{"2190-"21FF}
%\setmainjfont[
%  AddAltFont={Range=arrows, Font=STIX Two Math}
%]{Noto Sans CJK JP}



%\usepackage[style=ieee,citestyle=numeric-comp, maxbibnames=99, giveninits=false]{biblatex}
% Input centralized biblatex configuration
\input{lib/biblatex-and-friends}



% Use the chapter* style for the heading of the bibliography
\defbibheading{bibliography}[\bibname]{%
\section*{#1}%   skip \centering 
}
  
\addbibresource{references.bib}    
\input{kth/kthcolors}



\input{lib/defines}  % load some additional definitions to make writing more consistent

\usepackage[plainpages=false, unicode=true]{hyperref}
\input{lib/includes-after-hyperref}
\usepackage{orcidlink}  %% to provide ORCID icons and links

%\usepackage[acronym, style=super, toc=false, nonumberlist, nomain, nopostdot=true, notranslate]{glossaries-extra}
% In the document preamble where glossaries-extra is loaded:
\usepackage[
  acronym,
  toc=false,
  nonumberlist,
  nomain,
  nopostdot=true, 
  nohypertypes={main, acronym}, % disables link generation to targets that do not exist
  notranslate,
]{glossaries-extra}
%% In overleaf, if this file is in a folder and not the root
%% using \makeglossaries will not work, as the Overleaf latexmk
%% will not run the program to take the output.acn file and make the output.acr file
%\makeglossaries
% However, you can get TeX to do the work with the following:
\makenoidxglossaries
% For details of the Overleaf make process, see https://www.overleaf.com/learn/how-to/How_does_Overleaf_compile_my_project%3F
%% Potentially, one could set up a latexmkrc file in the folder

\input{lib/acronyms}                %load the acronyms file

% Include Glossary ---
% Align the text expansion of the glossary entries
\newglossarystyle{mylong}{%
  \setglossarystyle{long}%
  \renewenvironment{theglossary}%
     {\begin{longtable}[l]{@{}p{\dimexpr 2cm-\tabcolsep}p{0.8\hsize}}}% <-- change the value here
     {\end{longtable}}%
 }
 
% define a left-aligned table cell that is ragged right
\newcolumntype{L}[1]{>{\raggedright\let\newline\\\arraybackslash\hspace{0pt}}p{#1}}

\usepackage{cleveref}           %% Replace Section with a symbol
\usepackage{metalogo}   % for \LuaLaTeX logo

%% Conventions for todo notes:
% Informational
%% \generalExpl{Comments/directions/... in English}
\newcommand*{\generalExpl}[1]{\todo[inline]{#1}}                

% Language-specific information (currently in English or Swedish)
\newcommand*{\engExpl}[1]{\todo[inline, backgroundcolor=kth-lightgreen40]{#1}} %% \engExpl{English descriptions about formatting}
\newcommand*{\sweExpl}[1]{\todo[inline, backgroundcolor=kth-lightblue40]{\foreignlanguage{swedish}{#1}}}  %% % \sweExpl{Text på svenska}

% warnings
\newcommand*{\warningExpl}[1]{\todo[inline, backgroundcolor=kth-lightred40]{#1}} %% \warningExpl{warnings}


  \RequirePackage{fontspec}
  \defaultfontfeatures{Ligatures={TeX}} % This enables TeX style ligatures such as ---, '', ``, and so on




    \babelfont{rm}{TeX Gyre Termes}
    \DeclareFontShape{TU}{TeXGyreTermes(0)}{md}{n}{<->sub * TeXGyreTermes(0)/m/n}{}
    \DeclareFontShape{TU}{TeXGyreTermes(0)}{sb}{n}{<->ssub * TeXGyreTermes(0)/b/n}{}
    \DeclareFontShape{TU}{TeXGyreTermes(0)}{sb}{it}{<->ssub * TeXGyreTermes(0)/b/it}{}
  
    %\setsansfont{TeX Gyre Heros}   %% Helvetica like font
    \babelfont{sf}{TeX Gyre Heros}
    %\setmonofont[Ligatures={NoCommon}, Numbers={Lining,Monospaced}]{TeX Gyre Cursor}  %% Courier like font
    %% Turn off ligature for tt font
    \babelfont{tt}[Ligatures={ResetAll}]{TeX Gyre Cursor}
 \begin{comment}
    \babelfont[english]{rm}{TeX Gyre Termes}
    \babelfont[english]{sf}{TeX Gyre Heros}
    \babelfont[english]{tt}{TeX Gyre Cursor}
\end{comment}

 %  \setmathfont{TeX Gyre Termes Math} %% a math font
    \usepackage{mathtools}
 \usepackage[warnings-off={mathtools-colon,mathtools-overbracket}]{unicode-math}
  % The [version=bold, FakeBold=1.2] is to avoid a warning about the lack of a bold font
  %\setmathfont{TeX Gyre Pagella Math}[version=bold, FakeBold=1.2] %% a font for math
  \setmathfont{STIX Two Math}[version=normal]
  %\setmathfont{TeX Gyre Pagella Math}[version=normal]
  %\setmathfont{TeX Gyre Pagella Math}[version=bold, BoldFeatures={FakeBold=1.5}]
  % For both XeLaTeX and LuaLatex for getting access to unicode symbols
  %\newfontfamily\myfont[CharacterVariant=1]{NewCM10-Regular.otf}
  % STIX Project (Scientific and Technical Information Exchange)
  % STIX Two Math does not have bold face - so we fake it
  %\newfontfamily\mystixmathfont[BoldFeatures={FakeBold=1.5}, BoldItalicFeatures={FakeBold=1.5}]{STIX Two Math}
  \newfontfamily\mystixmathfont{STIX Two Math}
  % STIX Two Math does not have a bold font, but it has bold symbols with an without serifs - but you manually have to use them, unless you are in math mode - then you can use \symbf{}
  % and this will return the bold serif version of the character
  \DeclareFontShape{TU}{STIXTwoMath(0)}{b}{n}{<->ssub * STIXTwoMath(0)/m/n}{}
  \DeclareFontShape{TU}{STIXTwoMath(0)}{sb}{n}{<->ssub * STIXTwoMath(0)/m/n}{}
  \newfontfamily\mystixtextfont{STIX Two Text}
  
    % To access the Stylistic Set 1 <ss01> such as ℋ
    \newfontfamily\mystixmathfontSSa{STIX Two Math}[StylisticSet=1]

    % use english as a fallback when in other languages
    \babelprovide[import, onchar=ids fonts]{english}

    
  % for new KTH cover
  % Load the Figtree font as it is used for the new KTH graphical profile
  % 

    \newfontfamily{\FigtreeFont}[Ligatures=TeX,
        Path=./Figtree/static/,
        Extension = .ttf,
        UprightFont=*-Regular,
        BoldFont=*-Bold,
        BoldItalicFont=*-BoldItalic,
        ItalicFont=*-Italic,
        %FontFace={l}{n}{*-Light},
        %FontFace={l}{it}{*-LightItalic},
        FontFace={md}{n}{*-Medium},
        FontFace={md}{it}{*-MediumItalic},
        FontFace={sb}{n}{*-Semibold},
        FontFace={sb}{it}{*-SemiBoldItalic},
        %FontFace={k}{n}{*-Black},
        %FontFace={k}{it}{*-BlackItalic},
        %FontFace={eb}{n}{Font=*-ExtraBold},
        %FontFace={eb}{it}{Font=*-ExtraBoldItalic}
        ]{Figtree}


\newfontfamily\pageNumberFont{Figtree} %% set the font to use for page numbering

\newfontfamily{\NotoEmojiFont}[Ligatures=TeX,
    Path=./Noto_Emoji/static/,
    Extension = .ttf,
    UprightFont=*-Regular,
    BoldFont=*-Bold,
    FontFace={l}{n}{*-Light.ttf},
    FontFace={md}{n}{*-Medium},
    FontFace={sb}{n}{*-SemiBold},
    ]{NotoEmoji}



\begin{comment}
   % To set the abstract headings in Figtree, we redefine the abstract environment to look at the language being used and use the appropriate font, with the default being Figtree
   % The languages that are automatically introduced by Polyglossia or Babel have a name of the form xxxfont and xxxfontsf; where xxxfont is the serif font and xxxfontsf is the sans serif font.
   % This means that for each language that Figtree does not support, you have to define the sans serif and serif font to use.

      \babelprovide[import, onchar=ids fonts]{greek}
      \babelprovide[import, onchar=ids fonts]{russian}
      \babelprovide[import, onchar=ids fonts]{vietnamese}
      \babelfont[greek, russian, vietnamese]{rm}{Noto Serif}
      \babelfont[greek, russian, vietnamese]{sf}{Noto Sans}
      \babelfont[greek, russian, vietnamese]{tt}{Noto Mono}

  
    \babelprovide[import, onchar=ids fonts]{hindi}
    \babelfont[hindi]{rm}{Noto Serif Devanagari}
    \babelfont[hindi]{sf}{Noto Sans Devanagari}
    \babelfont[hindi]{tt}{Noto Sans Devanagari} % Noto Mono does not have the glyphs

    %\babelprovide[import, onchar=ids fonts]{russian}
    %\babelfont[russian]{rm}{Noto Serif}
    %\babelfont[russian]{sf}{Noto Sans}
    %\babelfont[russian]{tt}{Noto Mono}

    \babelprovide[import, onchar=ids fonts]{chinese-simplified}
    \babelfont[chinese-simplified]{rm}{Noto Serif CJK SC}
    \babelfont[chinese-simplified]{sf}{Noto Sans CJK SC}
    \babelfont[chinese-simplified]{tt}{Noto Sans Mono CJK SC}
    
    \babelfont[chinese-traditional]{rm}{Noto Serif CJK TC}
    \babelfont[chinese-traditional]{sf}{Noto Sans CJK TC}
    \babelfont[chinese-traditional]{tt}{Noto Sans Mono CJK TC}
  
    \babelprovide[import, onchar=ids fonts]{japanese}
    \babelfont[japanese]{rm}{Noto Serif CJK JP}
    \babelfont[japanese]{sf}{Noto Sans CJK JP}
    \babelfont[japanese]{tt}{Noto Sans Mono CJK JP}
  
  % If you are going to use Arabic

    \babelprovide[import, onchar=ids fonts]{arabic}
    \babelprovide[import, onchar=ids fonts]{centralkurdish}
    \babelfont[arabic, centralkurdish]{rm}{Noto Naskh Arabic}
    \babelfont[arabic, centralkurdish]{sf}{Noto Sans Arabic}
    \babelfont[arabic, centralkurdish]{tt}{Noto Sans Arabic}
    % If one really needs a monospaced font, one might try Kawkab Mono
    % However, it seems that it is a work in progress - see https://makkuk.com/kawkab-mono/ and https://github.com/aiaf/kawkab-mono/tree/master

    %\babelprovide[import, onchar=ids fonts]{centralkurdish}
    %\babelfont[centralkurdish]{rm}{Noto Naskh Arabic}
    %\babelfont[centralkurdish]{sf}{Noto Sans Arabic}
    %\babelfont[centralkurdish]{tt}{Noto Sans Arabic}

      
  % If you are going to use Hebrew or Yiddish
\babelprovide[import, onchar=ids fonts]{hebrew}
\babelprovide[import, onchar=ids fonts]{yiddish}
\babelfont[hebrew, yiddish]{rm}{Noto Serif Hebrew}
\babelfont[hebrew, yiddish]{sf}{Noto Sans Hebrew}
\babelfont[hebrew, yiddish]{tt}{Noto Sans Hebrew}


    %\babelprovide[import, onchar=ids fonts]{vietnamese}
    %\babelfont[vietnamese]{rm}{Noto Serif}
    %\babelfont[vietnamese]{sf}{Noto Sans}
    %\babelfont[vietnamese]{tt}{Noto Mono}

\end{comment}
  
    % The Overleaf TeX Live includes these fonts, so there is little you have to do!
    % The list of such fonts is at https://www.overleaf.com/learn/latex/Questions%2FWhich_OTF_or_TTF_fonts_are_supported_via_fontspec%3F
    %
\newfontfamily{\NotoSansJPMonoFont}[Ligatures=TeX,
    ]{Noto Sans Mono CJK JP Regular}


\newfontfamily{\NotoSansJPFont}[Ligatures=TeX,
    ]{Noto Sans CJK JP}


\newfontfamily{\NotoSansFont}[Ligatures=TeX,
Extension = .ttf,
    UprightFont = NotoSans-Regular,
    BoldFont = NotoSans-Bold,
    ItalicFont = NotoSans-Italic,
    BoldItalicFont = NotoSans-BoldItalic
    ]{NotoSansCustomA}
    
\newfontfamily{\NotoSerifFont}[Ligatures=TeX,
Extension = .ttf,
    UprightFont = NotoSerif-Regular,
    BoldFont = NotoSerif-Bold,
    ItalicFont = NotoSerif-Italic,
    BoldItalicFont = NotoSerif-BoldItalic
    ]{NotoSerifCustomA}

\newfontfamily{\DejaVuSansFont}[Ligatures=TeX,]{DejaVu Sans}

% A font to be used in minted listings
\newfontfamily{\ListingMono}{Noto Sans Mono}[Ligatures=ResetAll]

% fonts for various symbols
\newfontfamily{\NotoSansSymbolsFont}{Noto Sans Symbols}[Ligatures=ResetAll]
\newfontfamily{\NotoSansSymbolsTwoFont}{Noto Sans Symbols 2}[Ligatures=ResetAll]



% color symbols - this would require changes in the font used in the U+2600-U+26FF-Miscellaneous-Symbols.tex file
% for example for characters: ☕, ⚽, ⛄, and ✈
%\newfontfamily{\NotoColorEmojiFont}{Noto Color Emoji}[Renderer=Harfbuzz]

% Phonetic Extensions, U+1D00 - U+1D7F
\input{unicode_blocks/U+1D00-U+1D7F-Phonetic-Extensions}
% Superscripts and Subscripts, U+2070 - U+209F
\input{unicode_blocks/U+2070-U+209F-Superscripts-and-Subscripts}
% Arrows, U+2190 - U+21FF
\input{unicode_blocks/U+2190-U+21FF-Arrows}
% Miscellaneous Symbols, U+2600 - U+26FF
\input{unicode_blocks/U+2600-U+26FF-Miscellaneous-Symbols}

% Reclassify Arrows block (U+2190 - U+21FF) as Western characters (ALChar)
%\luaexec{
%  for c = 0x2190, 0x21FF do
%    luatexja.charrange.set_range_alphabetic(c)
%  end
%}
% Tell LuaTeX-ja to treat U+21AA as a regular Western character
%\ltjsetparameter{chartype={0x21AA, ALChar}}



% Set up the header and footer for plain style pages
\fancypagestyle{plain}{%
\fancyhf{}% clear all header and footer fields
\fancyfoot[C]{\pageNumberFont\small\selectfont\thepage}
\renewcommand{\headrulewidth}{0pt}%
\renewcommand{\footrulewidth}{0pt}%
}
% Set up the header and footer
\fancyhead{}
%\fancyhead[RO]{\sffamily\small\leftmark\thinspace|\thinspace\thepage}
%\fancyhead[LE]{\sffamily\small\thepage\thinspace|\thinspace\leftmark}
% Without conversion to uppercase
% Note that we need to switch to the familydefault for the page numbers
\fancyhead[RO]{{\sffamily\small\nouppercase\leftmark}{\pageNumberFont\small\selectfont \thinspace|\thinspace \thepage}}
\fancyhead[LE]{{\pageNumberFont\small\selectfont\thepage \thinspace|\thinspace}{\sffamily\small\nouppercase\leftmark}}
\fancyfoot{}
\renewcommand{\headrulewidth}{0pt}
\pagestyle{fancy}
\renewcommand{\sectionmark}[1]{%
\markboth{#1}{}}


\setlength{\headheight}{15pt}
\addtolength{\topmargin}{-3pt}

\input{lib/configure_listings}
% If you use both listing and lstlisting environments - the following makes them both use the same counter and the same formatting in the List of Listings
\makeatletter
\AtBeginDocument{\let\c@listing\c@lstlisting}
\AtBeginDocument{\let\l@listing\l@lstlisting}
\makeatother

\newcommand{\dname}[1]{\textbf{#1}}
\newcommand{\fname}[1]{\texttt{#1}}


\input{kth/Lua_timestamping}

% Input file generators before \documentclass
\input{lib/ext-eprint-dbx}

\input{lib/check_for_placeholder_references}

\title{README and notes about the template for programmers}

\author{Gerald Q. Maguire Jr.}
\date{October 2025}

\begin{document}
\pagenumbering{roman}
\glsresetall
\maketitle
\fancypagestyle{empty}{}
\warningExpl{\textbf{This document is a work in progress.}}

\begin{abstract}
This document describes the thesis template I developed for use at \gls{KTH} and explains why it is structured this way. This document provides programmers with information about the template, particularly how to use it with specific programs to support the thesis process for students, faculty, and administrators. The document is primarily intended for readers who are comfortable with programming or who need to maintain or modify the template or the programs I have written.

\warningExpl{{\mystixmathfont ☡} This is a very deep dive that is probably only of interest to a \textit{very} small number of readers.}
\end{abstract}
\cleardoublepage

%%%%%%%%%%%%%%%%%%%%%%%%%%%%%%%%%%%%%%%%%%%%%%%%%%%%%%%%%%%%%%%%%%%%%%
% Enable the following command the first time you compile this in a new session to give all of the time quanta to set up the font cache.
%% If the font cache is not built, you are likely to get a timeout.
%\end{document}

% Include Table of Contents (TOC) ---
\fancypagestyle{plain}{}
\renewcommand{\sectionmark}[1]{ \markboth{#1}{}}
\tableofcontents
\markboth{\contentsname}{}
\cleardoublepage
 
\setcounter{page}{1}
\pagenumbering{arabic}
\glsresetall
\section{Introduction}
\label{sec:IntroForProgrammers}

The focus in the first part of this document is on how to use the information in the thesis and captured by the \LaTeX{}  template and stored in the \fname{fordiva.json file} to do other tasks:
\begin{enumerate}
    \item The announcements and calendar events of thesis presentations or defenses. Moreover, thesis presentations and defenses are supposed to be public, so it would be useful and foster their visibility if it were easy to make such announcements (in both English and Swedish).
    
    \item Both the metadata and the thesis itself need to be entered into \foreignlanguage{swedish}{\gls{DiVA}}\footnote{There is a separate issue regarding whether the full-text is publicly available or not, but this is outside the scope of this document.}.
    
    \item Both English and Swedish titles will be reported along with the decision of the grading committee in \foreignlanguage{swedish}{\gls{LADOK}}.
\end{enumerate}


A high-level overview was given in the \fname{README\_supervisor\_notes.tex} file. This document covered \enquote{Accuracy and Economics}, \enquote{Make it simple from the start} by collecting the metadata, 
\enquote{Automating later steps in the degree project},
\enquote{Making an announcement},
\enquote{Making a cover and applying it},
\enquote{Can someone else use these programs?},
\enquote{\LaTeX{}~template in Overleaf}, and
\enquote{Put information into \LaTeX{}~template to generate a draft or final thesis}.

\Cref{sec:someEhnahncements} discusses some enhancements. Some readers may be interested in accessibility, see \Cref{sec:accessibility}. \Cref{sec:structureOfTheTemplate} discusses the structure of the template and the report. 
\Cref{sec:GeneratingLatexCommands} shows how one can generate \LaTeX{}  commands for use in the student’s thesis source file, specifically for the \texttt{custom\_configuration file.}
\Cref{sec:TheBiGPicture} presents the big picture of how the different data and programs interact.

\Cref{sec:copyrightOrCClicenseExaminer} presents how the copyright or Creative Commons License can be realized on the \texttt{bookinfo} page.
\Cref{sec:acronymsInForDIVAdata} describes how one can add the  \fname{acronyms.tex} file to the \gls{PDF} both as text and attach it to the \gls{PDF} file.
\Cref{sec:addingMetadataToPDFfile} presents the details of how to add additional metadata to the \gls{PDF} file, such that human readers and other tools can access this data. This is a rather expansive area, as there are those who would like to be able to mechanically extract data about the thesis and extract data from it (\eg citation data).
Finally, \Cref{sec:futurework} suggests some future work, such as how to better exploit unicode, better math support in \gls{DiVA}, and Cora-based \gls{DiVA}.

I hope that this document contains sufficient information to enable someone else to address any remaining problems (that will inevitably occur) and extend it in new ways to better support the whole thesis process and the various readers.

\section{Overall flow with an extension}
The approach taken in the overall work was to make data readily available via the template and collect \& output it in a 
“For DIVA” set of information at the end of the \gls{PDF} document and a \fname{fordiva.json} file. Additionally, we try to minimize the need to manually enter information by taking information from a Canvas course room (if one is relevant).

The overall flow is shown in \Cref{fig:LogicalOverviewOfWhatIsDone2}. Notice that this figure has been extended to include a JSON\_to\_Cora\_DiVA process that can use the Cora \gls{API} to directly create an entry in the new Cora-based \gls{DiVA}:
\tikzset{
    processBox/.style={rectangle, rounded corners, minimum width=3cm, minimum height=1cm,text centered, font=\sffamily, draw=black, fill=red!20},
    destinationBox/.style={rectangle, rounded corners, minimum width=3cm, minimum height=1cm,text centered, draw=black, fill=green!10},
    arrow/.style={thick,->,>=stealth}
}

\begin{figure}[!ht]
\resizebox{\textwidth}{!}{%
\begin{tikzpicture}
[align=left,node distance=2cm]


\node (latexFile) [tape,tape bend top=none,draw,font=\sffamily] {\LaTeX file};
\node (PDFfile) [tape,tape bend top=none,draw,font=\sffamily, above of=latexFile] {PDF file};
\node (DOCXFile) [tape,tape bend top=none,draw,font=\sffamily, below of=latexFile] {DOCX file};

\node (extractor) [processBox, right=1cm of latexFile] {Extractor};
\node (jsonFile) [tape,tape bend top=none,draw,font=\sffamily, right=0.5cm of extractor] {JSON file};
\node (start) [processBox, right=0.5cm of jsonFile] {JSON\_to\_calendar};

\node (jsontoMODS) [processBox, below of=start] {JSON\_to\_MODS};
\node (modsFile) [tape,tape bend top=none,draw,font=\sffamily, below of=jsontoMODS] {MODS file};
\node (diva) [destinationBox, right=1cm of modsFile] {Import into DiVA};


\node (jsontoCoraDiVA) [processBox, below= 3cm of modsFile] {JSON\_to\_Cora\_DiVA};
\node (DiVAentry) [destinationBox, right=1cm of jsontoCoraDiVA] {DiVA entry};

\node (calendarEvent) [destinationBox, right=1cm of start] {Canvas calendar event};
\node (announcement) [destinationBox,  above of =calendarEvent] {Canvas announcement};
\node (cortinaCalendarEvent) [destinationBox,  right=0.5cm of start, below of=calendarEvent] {KTH Cortina calendar event};
\draw [arrow] (latexFile) --  (extractor.west);
\draw [arrow] (PDFfile) --  (extractor.west);
\draw [arrow] (DOCXFile) --  (extractor.west);
\draw [arrow] (extractor) --  (jsonFile.west);
\draw [arrow] (jsonFile) --  (start.west);
\draw [arrow] (start.east) --  (announcement.west);
\draw [arrow] (start.east) -- (calendarEvent);
\draw [arrow] (start.east) -- (cortinaCalendarEvent.west);
\draw [arrow] (jsonFile) --  (jsontoMODS.west);
\draw [arrow] (jsontoMODS) --  (modsFile);
\draw [arrow] (modsFile) --  (diva);

\draw [arrow] (jsonFile) --  (jsontoCoraDiVA.north);
\draw [arrow] (jsontoCoraDiVA) --  (DiVAentry);
\end{tikzpicture}
}
\caption{From a \gls{PDF} file, a \LaTeX\;project, or \gls{DOCX} file -- generate an announcement, calendar entries, and a \gls{MODS} file for import into \gls{DiVA}}
  \label{fig:LogicalOverviewOfWhatIsDone2}
\end{figure}
\FloatBarrier

\warningExpl{☡  If you are still reading at this point, it is going to get deeper!}

\section{The big picture}
\label{sec:TheBiGPicture}

\Cref{tab:inforAndItsUses} shows the relationship between the information and the scripts that use it. The sections highlighted in yellow are optional. For example, only 1\textsuperscript{st} cycle theses will have a second author. While all theses will have two supervisors, some may have more (the scripts support up to 10).

Individuals can be associated with \gls{KTH}, in this case, they have a Local User Id and an organization element (with the implicit top level (L0) being \gls{KTH}). The organization levels of L1 == School and L2 == Department. It is possible to extend this further, but the programs currently only support L1 and L2. The Title and Alternative title fields are used to support the English and Swedish title (or visa versa) -– similarly, the abstracts and keywords must include English and Swedish versions; with additional support for abstracts and keywords in many languages. Where a school should be specified, one can use \texttt{\textbackslash schoolAcronym\{xxxx\}}.

The Canvas course room’s course\_id is often an argument to my scripts. One reason for this is to be able to use the information in the course room to translate a local user ID (\ie a kthid) to an integration ID (\ie a LADOK user ID) or to be able to look up details of an author or supervisor.

Note that for third-cycle degrees, the degree information is a little different as there are no program codes, but rather education codes. Also, there is a \texttt{\textbackslash degreeModifier\{Philosophy\}} macro to enable the user to specify Philosophy as a modifier to the degree name.

\begin{landscape}
\small{
\begin{longtable}{L{2cm}|L{4cm}|l|L{4cm}|L{4cm}|l l l l|}
\caption{Information and how it is used}\label{tab:inforAndItsUses}\\
    &  & & & \textbf{DOCX DocProperty} &  &  &  & \\
    \textbf{Top level} & \textbf{elements} & \textbf{subelements} & \textbf{\LaTeX{}~Macro} & \textbf{or fields} & \begin{rotate}{60}\textbf{calendar}\end{rotate} & \begin{rotate}{60}\textbf{cover}\end{rotate} & \begin{rotate}{60} \textbf{LADOK}\end{rotate} & \begin{rotate}{60}\textbf{MODS}\end{rotate}\\
        \endfirsthead
    \multicolumn{9}{c}%
{\tablename\ \thetable\ -- \textit{Continued from previous page}} \\
    &  &	&  & \textbf{DOCX DocProperty} &  &  &  & \\
    \textbf{Top level} & \textbf{elements} & \textbf{subelements}	& \textbf{\LaTeX~Macro} & \textbf{or fields} & \begin{rotate}{60}\textbf{calendar}\end{rotate} & \begin{rotate}{60}\textbf{cover}\end{rotate} & \begin{rotate}{60} \textbf{LADOK}\end{rotate} & \begin{rotate}{60}\textbf{MODS}\end{rotate}\\
\hline
\endhead
      \hline
Author1	\\
&	Last name &	&	\textbackslash authorsLastname\{\} &	Author1\_Last\_name & x &	x &	x &	x\\
&	First name & &	\textbackslash authorsFirstname\{\} &	Author1\_First\_name & x & x & x & x\\
&	Local User Id &	& \textbackslash kthid\{\} &	Author1\_Local User Id &	& 		& x &	x \\
&	E-mail	& &	\textbackslash email\{\} &	Author1\_E-mail	&	& & &		x \\
 & organisation & \\
 &  & L1 & \textbackslash authorsSchool\{\} & Author1\_organization\_L1 &  &  &  & x\\
 &  &  &  & Author1\_organization\_L2 &  &  &  & \\
 & Other organisation &  &  & Author1\_Other\_organisation &  &  &  & \\
       \hline
\rowcolor{yellow}Author2 & \multicolumn{8}{l}{\cellcolor{yellow}}\\
\rowcolor{yellow} & Last name &  & \textbackslash secondAuthorsLastname\{\} & Author2\_Last\_name & x & x & x & x\\
\rowcolor{yellow} & First name &  & \textbackslash secondAuthorsFirstname\{\} & Author2\_First\_name & x & x & x & x\\
\rowcolor{yellow} & Local User Id &  & \textbackslash secondkthid\{\} & Author2\_Local User Id &  &  & x & x\\
\rowcolor{yellow} & E-mail &  & \textbackslash secondemail\{\} & Author2\_E-mail &  &  &  & x\\
\rowcolor{yellow} & organisation & \multicolumn{7}{l}{\cellcolor{yellow}}\\ 
\rowcolor{yellow} &  & L1 & \textbackslash secondAuthorsSchool{ } & Author2\_organization\_L1 &  &  &  & x\\
\rowcolor{yellow} &  &  &  & Author2\_organization\_L2 &  &  &  & \\
\rowcolor{yellow} & Other organisation &  &  & Author2\_Other\_organisation &  &  &  & \\
       \hline
Cycle &  &  & \textbackslash courseCycle\{\} & Cycle & x & x &  & \\
Course code &  &  & \textbackslash courseCode\{\} & Course\_code & i & i &  &  \\
Credits &  &  & \textbackslash courseCredits\{\} & Credits &  &  &  & \\
      \hline
Degree1 & \\
 & Educational program &  & Derived from programcode & Educational program &  & x &  & x \\
 & programcode &  & \textbackslash programcode\{\} & programcode &  & i &  & i \\
 & Degree &  & \textbackslash degreeName\{\} & Degree &  & x &  & x \\
 & subjectArea &  & \textbackslash subjectArea{ } & subjectArea &  & x &  & x \\
\rowcolor{yellow}Degree2 & \multicolumn{8}{l}{\cellcolor{yellow}}\\
\rowcolor{yellow}& Educational program &  & Derived from secondProgramcode & second\_Educational program &  & x &  & x \\
\rowcolor{yellow}& programcode &  & \textbackslash secondProgramcode\{\} & second\_Programcode &  & i &  & i \\
\rowcolor{yellow}& Degree &  & \textbackslash secondDegreeName\{\} & second\_Degree &  & x &  & x \\
\rowcolor{yellow}& subjectArea &  & \textbackslash secondSubjectArea\{\} & second\_subjectArea &  & x &  & x \\
       \hline
Title & \\
 & Main title &  & \textbackslash title\{\} & (exiting document property) & x & x & x & x \\
 & Subtitle &  & \textbackslash subtitle\{\} & Subtitle & x & x & o & x \\
 & Language &  & Derived from option to document class &  & x &  &  & x \\
Alternative title & \\
 & Main title &  & \textbackslash alttitle\{\} & Alternative\_main\_title & x & x & x & x \\
 & Subtitle &  & \textbackslash altsubtitle\{\} & Alternative\_subtitle & x & x & o & x \\
 & Language &  & Derived from option to documentclass &  & x &  &  & x \\
      \hline
Supervisor1 & \\
 & Last name &  & \textbackslash supervisorAsLastname\{\} & Supervisor1\_Last\_name & x &  &  & x \\
 & First name &  & \textbackslash supervisorAsFirstname\{\} & Supervisor1\_First\_name & x &  &  & x \\
 & Local User Id &  & \textbackslash supervisorAsKTHID\{\} & Supervisor1\_Local User Id &  &  &  & x \\
 & E-mail &  & \textbackslash supervisorAsEmail\{\} & Supervisor1\_E-mail &  &  &  & x \\
 & organisation & \\
 &  & L1 & \textbackslash supervisorAsSchool\{\} & Supervisor1\_organization\_L1 &  &  &  & x \\
 &  & L2 & supervisorAsDepartment\{\} & Supervisor1\_organization\_L2 &  &  &  & o \\
 & Other organisation &  & supervisorAsOrganization\{\} & Supervisor1\_Other\_organisation &  &  &  &  \\
\rowcolor{yellow}Supervisor2 &  \multicolumn{8}{l}{\cellcolor{yellow}}\\
\rowcolor{yellow} & Last name &  & \textbackslash supervisorBsLastname{ } & Supervisor2\_Last\_name & x &  &  & x \\
\rowcolor{yellow} & First name &  & \textbackslash supervisorBsFirstname\{\} & Supervisor2\_First\_name & x &  &  & x \\
\rowcolor{yellow} & Local User Id &  & \textbackslash supervisorBsKTHID\{\} & Supervisor2\_Local User Id &  &  &  & x \\
\rowcolor{yellow} & E-mail &  & \textbackslash supervisorBsEmail\{\} & Supervisor2\_E-mail &  &  &  & x \\
\rowcolor{yellow} & organisation & \multicolumn{7}{l}{\cellcolor{yellow}}\\
\rowcolor{yellow} &  & L1 & \textbackslash supervisorBsSchool\{\} & Supervisor2\_organization\_L1 &  &  &  & x \\
\rowcolor{yellow} &  & L2 & \textbackslash supervisorBsDepartment\{\} & Supervisor2\_organization\_L2 &  &  &  & o \\
\rowcolor{yellow} & Other organisation &  & \textbackslash supervisorBsOrganization\{\} & Supervisor2\_Other\_organisation &  &  &  & \\
\rowcolor{yellow}Supervisor3 & \multicolumn{8}{l}{\cellcolor{yellow}} \\
\rowcolor{yellow} & Last name &  & \textbackslash supervisorCsLastname\{\} & Supervisor3\_Last\_name & x &  &  & x \\
\rowcolor{yellow} & First name &  & \textbackslash supervisorCsFirstname\{\} & Supervisor3\_First\_name & x &  &  & x \\
\rowcolor{yellow} & Local User Id &  & \textbackslash supervisorCsKTHID\{\} &  &  &  &  & \\ 
\rowcolor{yellow} & E-mail &  & \textbackslash supervisorCsEmail\{\} & Supervisor3\_E-mail &  &  &  & o \\
\rowcolor{yellow} & organisation & \multicolumn{7}{l}{\cellcolor{yellow}}\\ 
\rowcolor{yellow} &  & L1 & \textbackslash supervisorCsSchool\{\} & Supervisor3\_organization\_L1 &  &  &  & x \\
\rowcolor{yellow} &  & L2 & \textbackslash supervisorCsDepartment\{\} & Supervisor3\_organization\_L2 &  &  &  & o \\
\rowcolor{yellow} & Other organisation &  & \textbackslash supervisorCsOrganization\{\} & Supervisor3\_Other\_organisation & x &  &  &  x \\
 \hline
\rowcolor{yellow}Cooperation & \multicolumn{8}{l}{\cellcolor{yellow}}\\
\rowcolor{yellow} & Partner\_name &  & \textbackslash hostcompany\{\}        or \textbackslash hostorganization\{\} & Cooperation\_Partner\_name &  &  &  & o \\
  \hline
National Subject Categories &  &  & \textbackslash nationalsubjectcategories\{\} & National Subject Categories &  &  &  & x \\
 \hline
Other information & \\
 & Year &  & \textbackslash the\textbackslash year & (derived from document property Date) &  &  &  & x \\
 & Number of pages &  & (derived from labels: pg:lastPageofPreface and pg:lastPageofMainmatter) & (derived from bookmarks: lastPageofPreface and lastPageofMainmatter) &  &  &  & x \\
Series & \\
 & Title of series &  & \textbackslash trita{TRITA-EECS-EX}{2021:00} & Series\_name &  & x &  & x \\
 & No. in series &  &  & Number\_in\_series &  & x &  & x \\
Opponents &  &  & \textbackslash opponentsNames\{\} & Opponents\_Name & x &  &  & \\
Presentation & \\
 & Date &  & \textbackslash presentationDateAndTimeISO\{\} & Presentation\_Date & x &  &  & x \\
 & Language &  & \textbackslash presentationLanguage{ } & Presentation\_Language & x &  &  & x \\
 & Room &  & \textbackslash presentationRoom\{\} & Presentation\_Room & x &  &  & x \\
 & Address &  & \textbackslash presentationAddress{ } & Presentation\_Address & x &  &  & x \\
 & City &  & \textbackslash presentationCity{ } & Presentation\_City & x &  &  & x \\
\hline
Number of lang instances &  &  &  & [Manually entered as 2, 3, 4, …] &  &  &  & \\
abstracts & \\
 & eng &  & (saved in Lua table) & (derived from bookmark: EnglishAbstract and then text inserted) & x &  &  & x \\
 & swe &  &  & (derived from bookmark: SwedishAbstract) & x &  &  & x \\
 & \cellcolor{yellow}fre &\cellcolor{yellow}  & & \multicolumn{4}{l}{\cellcolor{yellow}} & \cellcolor{yellow}o\\
 & \cellcolor{yellow}… & \cellcolor{yellow} & & \multicolumn{4}{l}{\cellcolor{yellow}} & \cellcolor{yellow}o\\
keywords & \\
 & eng &  & (saved in a Lua tables) & (derived from bookmark: EnglishKeywords) & x &  &  & x \\
 & swe &  &  & (derived from bookmark: SwedishKeywords) & x &  &  & x \\
 & \cellcolor{yellow}fre & \cellcolor{yellow} & & \multicolumn{4}{l}{\cellcolor{yellow}} & \cellcolor{yellow}o \\
 & \cellcolor{yellow}… & \cellcolor{yellow} & & \multicolumn{4}{l}{\cellcolor{yellow}} & \cellcolor{yellow}o\\
\hline
\end{longtable}
}
\end{landscape}


\section{Use of Lua and Babel}
One of the major changes introduced in this template is the use of Lua and Babel. Lua is a simple programming language. This language is available to the user of \LuaLaTeX{}. Lua's tables offer a very powerful way to store structured data. We use a Lua package (\fname{dkjson.lua}) to write a Lua table containing all of the metadata to a \fname{fordiva.json} file. The format of this file is \gls{JSON}. Not only is this file generated but it is both attached to the PDF file and included in a listing on the \enquote{For DiVA} page.

The primary use of the \fname{fordiva.json} file is to facilitate reporting the thesis in \gls{DiVA}. However, a secondary use of this information is to potentially automate many of the steps in announcing the oral presentation (as was described in the section \enquote{Making an announcement} in the file \fname{README\_supervisor\_notes.tex}).
For example, the student's organization or main supervisor's organization is used to determine which local part of the \gls{KTH} calendar a calendar announcement should appear in, as the Cortina calendar is divided by the school and then by department.

\warningExpl{Note that the department name must be in Swedish for Cortina. Ideally, this is something that a program using the JSON data should handle for the user.}

In addition to the English and Swedish abstracts and keywords information, there are a number of skeleton files in the folder \texttt{Additional\_Abstracts}. While not an exhaustive list of possibilities, they represent the languages for which I found several abstracts in DiVA\footnote{Except for Estonian, which had an entry but no longer has an entry.} Since 2000, the language code \enquote{scc} (for Serbo-Croatian) is considered deprecated and has been replaced by \enquote{srp} for Serbian and \enquote{hrv} for Croatian (for which I have not implemented a skeleton). Otherwise, there is a skeleton abstract for all of the language codes used in one or more abstracts in this language in \gls{DiVA}.

The \texttt{\textbackslash foreignlanguage\{xxxx\}\{xxxxxxx\}} command can be used to indicate that the string xxxxxxx is in language xxxx and the use of the \texttt{\textbackslash selectlanguage\{xxxx\}} command to switch to the language xxxx. The first command makes it easy to include text in a language that is not the main language of the thesis at any point in the thesis. The second command can be used to switch to a specific language at a certain point in the text. This second command is used to facilitate the inclusion of abstracts and keywords in various languages.

The \texttt{babel} package is used to handle different languages and automatically switch to the appropriate font for each language. 
If one does \textbf{not} use babel configured for the languages that you use, you will find messages in the \fname{output.log} file, such as those shown in \Cref{lst:missingCharactersFromLogExample}. While it is possible to work around this on a character-by-character basis by using suitable fonts and adding mappings, such as shown in \Cref{lst:newunicodeForAbstracts}, a better approach is described in \Cref{sec:BetterHandlingOfFonts}.


\begin{lstlisting}[style=latexExampleForAuthors,  escapechar=|, 
basicstyle=\footnotesize,
caption={Examples of missing characters from the output.log file},
label=lst:missingCharactersFromLogExample]
\end{lstlisting}
\begin{Verbatim}[breaklines=true]
Missing character: There is no 中 (U+4E2D) in font TeX Gyre Cursor/OT:script=latn;language=dflt;-liga;+lnum;+tnum; mapping=tex-text;!
Missing character: There is no 文 (U+6587) in font TeX Gyre Cursor/OT:script=latn;language=dflt;-liga;+lnum;+tnum; mapping=tex-text;!
Missing character: There is no 摘 (U+6458) in font TeX Gyre Cursor/OT:script=latn;language=dflt;-liga;+lnum;+tnum; mapping=tex-text;!
Missing character: There is no 要 (U+8981) in font TeX Gyre Cursor/OT:script=latn;language=dflt;-liga;+lnum;+tnum; mapping=tex-text;!
Missing character: There is no 个 (U+4E2A) in font TeX Gyre Heros/OT:script=latn;language=dflt;mapping=tex-text;!
Missing character: There is no 关 (U+5173) in font TeX Gyre Heros/OT:script=latn;language=dflt;mapping=tex-text;!
Missing character: There is no 键 (U+952E) in font TeX Gyre Heros/OT:script=latn;language=dflt;mapping=tex-text;!
Missing character: There is no 词 (U+8BCD) in font TeX Gyre Heros/OT:script=latn;language=dflt;mapping=tex-text;!
\end{Verbatim}


\begin{lstlisting}[style=latexExampleForAuthors,  escapechar=|, 
caption={Examples of new unicode character definitions},
label=lst:newunicodeForAbstracts]
\end{lstlisting}
\begin{Verbatim}[breaklines=true]
\newunicodechar{个}{\iffontchar\font`个 个\else{\cjkfonttt 个}\fi} % U+4E2A - CJK Unified Ideograph-4E2A
\newunicodechar{中}{\iffontchar\font`中 中\else{\cjkfonttt 中}\fi} % U+4E2D - CJK Unified Ideograph-4E2D
\newunicodechar{关}{\iffontchar\font`关 关\else{\cjkfonttt 关}\fi} % U+5173 - CJK Unified Ideograph-5173
\newunicodechar{摘}{\iffontchar\font`摘 摘\else{\cjkfonttt 摘}\fi} % U+6458 - CJK Unified Ideograph-6458
\newunicodechar{文}{\iffontchar\font`文 文\else{\cjkfonttt 文}\fi} % U+6587 - CJK Unified Ideograph-6587
\newunicodechar{要}{\iffontchar\font`要 要\else{\cjkfonttt 要}\fi} % U+8981 - CJK Unified Ideograph-8981
\newunicodechar{词}{\iffontchar\font`词 词\else{\cjkfonttt 词}\fi} % U+8BCD - CJK Unified Ideograph-8BCD
\newunicodechar{键}{\iffontchar\font`键 键\else{\cjkfonttt 键}\fi} % U+952E - CJK Unified Ideograph-952E
\end{Verbatim}

\FloatBarrier


If the user uses characters that require fonts other than \texttt{TeX Gyre Termes} and has not configured babel with a font for this language, the contents of the abstracts in these languages may not appear on the \enquote{For DiVA} page but will be in the \fname{fordiva.json} file. This will occur because the file simply contains the UTF-8 encoded unicode for the abstracts.

\subsection{Better handling of fonts}
\label{sec:BetterHandlingOfFonts}

Manually or programmatically adding all of the missing characters to the relevant \texttt{unicode\_blocks} file was used to work with polyglossia and \XeLaTeX{}. However, it took a lot of work.

Therefore, a decision was made to switch to \LuaLaTeX{} and babel. This combination has well-integrated \texttt{fontspec} support for automatically switching language and fonts. \Cref{lst:BabelHindiExample} shows how easy it is to set up a font that can be used when setting content written in Hindi, \ie babel language hindi. This font will automatically be invoked when you issue a \texttt{\textbackslash selectlenguage\{hindi\}} \textbf{and} when you use one of the unicode characters with code points in the Devanagari range, \ie U+0900 to U+097F.

\begin{lstlisting}[language={[LaTeX]TeX}, caption={Establishing fonts to be used when setting content in Hindi}, label=lst:BabelHindiExample] 
\babelprovide[import, onchar=ids fonts]{hindi}
\babelfont[hindi]{rm}{Noto Serif Devanagari}
\babelfont[hindi]{sf}{Noto Sans Devanagari}
\end{lstlisting}

Similar functionality is utilized by the template with the babel languages: arabic, centralkurdish, chinese-simplified, chinese-traditional greek, hebrew, japanese, russian, and vietnamese.

One of the advantages of switching to babel is that it allows lazy loading of languages and lazy loading of fonts -- thus unless the language is used or the font is used, they do not have to be loaded.

Unfortunately, \texttt{luaotfload} program must still resolve the font name into a file to load the font from, then it transforms this font into Lua code and compiles it into bytecode (that later can be loaded -- if the font is used). The time required to do all of this can be considerable and may lead to timeouts. For this reason, I have introduced \fname{warmup.tex} that when compiled will cause the \texttt{luaotfload} program to load the fonts used by \fname{examplethesis.tex} into the working font cache. Addition, it will store a copy of these font files into a working directory. By following the directions in \fname{Saving\_and\_restoring\_font\_cache.tex} you can download the saved fonts and then upload them for use by the document \fname{restore\_font\_cache.tex} that when compiled will copy the uploaded saved fonts to the working font cache.

\subsubsection{Babel ranges}

The file \texttt{babel-data-bidi.lua} \footnote{Perhaps under \texttt{/usr/share/texmf/tex/generic/babel/}} defines the Babel table. This is shown in \Cref{lst:BabelTableandRanges}. The format of the entries is the starting codepoint of a range, the ending codepoint of a range (inclusive), and the direction. The direction is defined according to \enquote{Table 4. Bidirectional Character Types} in \textit{Unicode® Standard Annex \#9: The Bidirectional Algorithm}~\cite{UAX9}. Accordingly, `r' is right-to-left, `l' is is left-to-right (the default), `al' is right-to-left Arabic, and `on' is other neutrals.

\begin{lstlisting}[language={[5.2]Lua}, caption={Babel table and Babel-ranges supplemented with additional comments about Unicode blocks}, label=lst:BabelTableandRanges]
Babel = Babel or {}

Babel.ranges={
 {0x0590, 0x05FF, 'r'}, -- Hebrew, U+0590 - U+05FF
 {0x0600, 0x07BF, 'al'}, -- Arabic, U+0600 - U+06FF;
                         -- Syriac, U+0700 - U+074F;
                         -- Arabic Supplement, U+0750 - U+077F;
                         -- Thaana, U+0780 - U+07BF
 {0x07C0, 0x085F, 'r'}, -- NKo, U+07C0 - U+07FF;
                        -- Samaritan, U+0800 - U+083F;
                        -- Mandaic, U+0840 - U+085F
 {0x0860, 0x086F, 'al'}, -- Syriac Supplement, U+0860 - U+086F
 {0x08A0, 0x08FF, 'al'}, -- Arabic Extended-A, U+08A0 - U+08FF
 {0xE000, 0xF8FF, 'on'}, -- PUA -- Private Use Area, U+E000 - U+F8FF
 {0xFB1D, 0xFB4F, 'r'},  -- Alphabetic Presentation Forms, U+FB00 - U+FB4F
 {0xFB50, 0xFDFF, 'al'}, -- Arabic Presentation Forms-A, U+FB50 - U+FDFF
 {0xFE70, 0xFEFF, 'al'}, -- Arabic Presentation Forms-B, U+FE70 - U+FEFF
 {0x10800, 0x10C4F, 'r'}, -- Cypriot Syllabary, U+10800 - U+1083F;
                          -- Imperial Aramaic, U+10840 - U+1085F;
                          -- Palmyrene, U+10860 - U+1087F;
                          -- Nabataean, U+10880 - U+108AF; 
                        -- hole U+108B0 - 108DF;
                         -- Hatran, U+108E0 - U+108FF;
                         -- Phoenician, U+10900 - U+1091F;
                         -- Lydian, U+10920 - U+1093F;
                         -- Meroitic Hieroglyphs, U+10980 - U+1099F;
                         -- Meroitic Cursive, U+109A0 - U+109FF;
                         -- Kharoshthi, U+10A00 - U+10A5F;
                         -- Old South Arabian, U+10A60 - U+10A7F;
                         -- Old North Arabian, U+10A80 - U+10A9F;
                         -- Manichaean, U+10AC0 - U+10AFF;
                         -- Avestan, U+10B00 - U+10B3F;
                         -- Inscriptional Parthian, U+10B40 - U+10B5F;
                         -- Inscriptional Pahlavi, U+10B60 - U+10B7F;
                         -- Psalter Pahlavi, U+10B80 - U+10BAF;
                         -- Old Turkic, U+10C00 - U+10C4F
 {0x1E800, 0x1E8DF, 'r'}, -- Mende Kikakui, U+1E800 - U+1E8DF
 {0x1E900, 0x1E95F, 'r'}, -- Adlam, U+1E900 - U+1E95F
 {0x1F300, 0x1F9FF, 'on'}, -- Miscellaneous Symbols and Pictographs, U+1F300 - U+1F5FF;
                           -- Emoticons, U+1F600 - U+1F64F;
                           -- Ornamental Dingbats, U+1F650 - U+1F67F;
                           -- Transport and Map Symbols, U+1F680 - U+1F6FF;
                           -- Alchemical Symbols, U+1F700 - U+1F77F;
                           -- Alchemical Symbols, U+1F700 - U+1F77F;
                           -- Supplemental Arrows-C, U+1F800 - U+1F8FF;
                           -- Supplemental Symbols and Pictographs, U+1F900 - U+1F9FF:
 {0xF0000, 0xFFFFD, 'on'}, -- PUA -- Supplementary Private Use Area-A, U+F0000 - U+FFFFF; does not include U+FFFFE and U+FFFFF 
 {0x100000, 0x10FFFD, 'on'} -- PUA -- Supplementary Private Use Area-B, U+100000 - U+10FFFF; does not include U+10FFFE or U+10FFFF
}
\end{lstlisting}

\Cref{lst:BabelTableChracters} shows part of the \texttt{Babel.characters} table. The direction table (with \texttt{d} in it) has additional character types:
`an' for Arabic number,
`b' for paragraph separator, 
`bn' for boundary neutral,
`cs' for common number separator,
`en'  for European number,
`es'  for European number separator,
`et' for European number terminator,
`fsi' for first strong isolate,
`lre' for left-to-right embedding,
`lri' for left-to-right isolate,
`lro' for left-to-right override,
`nsm' for nonspace mark,
`pdf' for pop directional format,
`pdi' for pop directional isolate,
`rle' for right-left embedding,
`rli' for right-left isolate,
`rlo' for right-left override,
`s' for segment separator, and
`ws' for white space.
\clearpage
\begin{lstlisting}[language={[5.2]Lua},  escapechar=£, caption={(part of) Babel table of characters}, label=lst:BabelTableChracters]
Babel.characters={
 [0x0]={d='bn'},
 £\quad$\vdots$£
 [0x8]={d='bn'},
 [0x9]={d='s'},
 [0xA]={d='b'},
 [0xB]={d='s'},
 [0xC]={d='ws'},
 [0xD]={d='b'},
 [0xE]={d='bn'},
 £\quad$\vdots$£
 [0x1B]={d='bn'},
 [0x1C]={d='b'},
 [0x1D]={d='b'},
 [0x1E]={d='b'},
 [0x1F]={d='s'},
 [0x20]={d='ws'},
 [0x21]={d='on'},
 [0x22]={d='on'},
 [0x23]={d='et'},
 [0x24]={d='et'},
 [0x25]={d='et'},
 [0x26]={d='on'},
 [0x27]={d='on'},
 [0x28]={d='on', m=0x29},
 [0x29]={d='on', m=0x28},
 [0x2A]={d='on'},
 [0x2B]={d='es'},
 [0x2C]={d='cs'},
 [0x2D]={d='es'},
 [0x2E]={d='cs'},
 [0x2F]={d='cs'},
 [0x30]={d='en'},
 £\quad$\vdots$£
 [0x39]={d='en'},
 [0x3A]={d='cs'},
 [0x3B]={d='on'},
 [0x3C]={d='on', m=0x3E},
 [0x3D]={d='on'},
 [0x3E]={d='on', m=0x3C},
 [0x3F]={d='on'},
 [0x40]={d='on'},
 [0x5B]={d='on', m=0x5D},
 [0x5C]={d='on'},
 [0x5D]={d='on', m=0x5B},
 [0x5E]={d='on'},
 [0x5F]={d='on'},
 [0x60]={d='on'},
 [0x7B]={d='on', m=0x7D},
 [0x7C]={d='on'},
 [0x7D]={d='on', m=0x7B},
 [0x7E]={d='on'},
 [0x7F]={d='bn'},
 £\quad$\vdots$£
 [0x84]={d='bn'},
 [0x85]={d='b'},
 [0x86]={d='bn'},
 £\quad$\vdots$£
 [0x9F]={d='bn'},
 [0xA0]={d='cs'},
 [0xA1]={d='on'},
 [0xA2]={d='et'},
 [0xA3]={d='et'},
 [0xA4]={d='et'},
 [0xA5]={d='et'},
 [0xA6]={d='on'},
 [0xA7]={d='on'},
 [0xA8]={d='on'},
 [0xA9]={d='on'},
 [0xAB]={d='on', m=0xBB},
 [0xAC]={d='on'},
 [0xAD]={d='bn'},
 [0xAE]={d='on'},
 [0xAF]={d='on'},
 [0xB0]={d='et'},
 [0xB1]={d='et'},
 [0xB2]={d='en'},
 [0xB3]={d='en'},
 [0xB4]={d='on'},
 [0xB6]={d='on'},
 [0xB7]={d='on'},
 [0xB8]={d='on'},
 [0xB9]={d='en'},
 [0xBB]={d='on', m=0xAB},
 [0xBC]={d='on'},
 [0xBD]={d='on'},
 [0xBE]={d='on'},
 [0xBF]={d='on'},
 [0xD7]={d='on'},
 [0xF7]={d='on'},
 [0x2B9]={d='on'},
 [0x2BA]={d='on'},
 [0x2C2]={d='on'},
 £\quad$\vdots$£
\end{lstlisting}

Entries of the form shown in \Cref{lst:BabelTableChractersMachedDelimiters} define \textbf{mirrored} characters.

\begin{lstlisting}[language={[5.2]Lua},  escapechar=£, caption={Some of the matching delimiters in Babel table of characters (with added comments)}, label=lst:BabelTableChractersMachedDelimiters]
 [0x28]={d='on', m=0x29}, -- Left Parenthesis
 [0x29]={d='on', m=0x28}, -- Right Parenthesis
 [0x3C]={d='on', m=0x3E}, -- Less-Than Sign
 [0x3E]={d='on', m=0x3C}, -- Greater-Than Sign
 [0x5B]={d='on', m=0x5D}, -- Left Square Bracket
 [0x5D]={d='on', m=0x5B}, -- Right Square Bracket
 [0x7B]={d='on', m=0x7D}, -- Left Curly Bracket
 [0x7D]={d='on', m=0x7B}, -- Right Curly Bracket
 [0xAB]={d='on', m=0xBB}, -- Left-Pointing Double Angle Quotation Mark
 [0xBB]={d='on', m=0xAB}, -- Right-Pointing Double Angle Quotation Mark
 [0xF3A]={d='on', m=0xF3B}, -- Tibetan Mark Gug Rtags Gyon
 [0xF3B]={d='on', m=0xF3A}, -- Tibetan Mark Gug Rtags Gyas
 [0xF3C]={d='on', m=0xF3D}, -- Tibetan Mark Ang Khang Gyon
 [0xF3D]={d='on', m=0xF3C}, -- Tibetan Mark Ang Khang Gyas
 [0x169B]={d='on', m=0x169C}, -- Ogham Feather Mark
 [0x169C]={d='on', m=0x169B}, -- Ogham Reversed Feather Mark
 £\quad$\vdots$£
\end{lstlisting}

\Needspace*{14\baselineskip}
The file \texttt{luababel.def} defines \texttt{Babel.script\_blocks}. These are shown in \Cref{lst:BabelTableScriptBlocks}.

\begin{lstlisting}[language={[5.2]Lua},  escapechar=£, caption={Babel table script\_blocks}, label=lst:BabelTableScriptBlocks]
Babel.script_blocks = {
  ['dflt'] = {},
  ['Arab'] = {{0x0600, 0x06FF}, {0x08A0, 0x08FF}, {0x0750, 0x077F},
              {0xFE70, 0xFEFF}, {0xFB50, 0xFDFF}, {0x1EE00, 0x1EEFF}},
  ['Armn'] = {{0x0530, 0x058F}},
  ['Beng'] = {{0x0980, 0x09FF}},
  ['Cher'] = {{0x13A0, 0x13FF}, {0xAB70, 0xABBF}},
  ['Copt'] = {{0x03E2, 0x03EF}, {0x2C80, 0x2CFF}, {0x102E0, 0x102FF}},
  ['Cyrl'] = {{0x0400, 0x04FF}, {0x0500, 0x052F}, {0x1C80, 0x1C8F},
              {0x2DE0, 0x2DFF}, {0xA640, 0xA69F}},
  ['Deva'] = {{0x0900, 0x097F}, {0xA8E0, 0xA8FF}},
  ['Ethi'] = {{0x1200, 0x137F}, {0x1380, 0x139F}, {0x2D80, 0x2DDF},
              {0xAB00, 0xAB2F}},
  ['Geor'] = {{0x10A0, 0x10FF}, {0x2D00, 0x2D2F}},
  % Don't follow strictly Unicode, which places some Coptic letters in
  % the 'Greek and Coptic' block
  ['Grek'] = {{0x0370, 0x03E1}, {0x03F0, 0x03FF}, {0x1F00, 0x1FFF}},
  ['Hans'] = {{0x2E80, 0x2EFF}, {0x3000, 0x303F}, {0x31C0, 0x31EF},
              {0x3300, 0x33FF}, {0x3400, 0x4DBF}, {0x4E00, 0x9FFF},
              {0xF900, 0xFAFF}, {0xFE30, 0xFE4F}, {0xFF00, 0xFFEF},
              {0x20000, 0x2A6DF}, {0x2A700, 0x2B73F},
              {0x2B740, 0x2B81F}, {0x2B820, 0x2CEAF},
              {0x2CEB0, 0x2EBEF}, {0x2F800, 0x2FA1F}},
  ['Hebr'] = {{0x0590, 0x05FF}},
  ['Jpan'] = {{0x3000, 0x303F}, {0x3040, 0x309F}, {0x30A0, 0x30FF},
              {0x4E00, 0x9FAF}, {0xFF00, 0xFFEF}},
  ['Khmr'] = {{0x1780, 0x17FF}, {0x19E0, 0x19FF}},
  ['Knda'] = {{0x0C80, 0x0CFF}},
  ['Kore'] = {{0x1100, 0x11FF}, {0x3000, 0x303F}, {0x3130, 0x318F},
              {0x4E00, 0x9FAF}, {0xA960, 0xA97F}, {0xAC00, 0xD7AF},
              {0xD7B0, 0xD7FF}, {0xFF00, 0xFFEF}},
  ['Laoo'] = {{0x0E80, 0x0EFF}},
  ['Latn'] = {{0x0000, 0x007F}, {0x0080, 0x00FF}, {0x0100, 0x017F},
              {0x0180, 0x024F}, {0x1E00, 0x1EFF}, {0x2C60, 0x2C7F},
              {0xA720, 0xA7FF}, {0xAB30, 0xAB6F}},
  ['Mahj'] = {{0x11150, 0x1117F}},
  ['Mlym'] = {{0x0D00, 0x0D7F}},
  ['Mymr'] = {{0x1000, 0x109F}, {0xAA60, 0xAA7F}, {0xA9E0, 0xA9FF}},
  ['Orya'] = {{0x0B00, 0x0B7F}},
  ['Sinh'] = {{0x0D80, 0x0DFF}, {0x111E0, 0x111FF}},
  ['Syrc'] = {{0x0700, 0x074F}, {0x0860, 0x086F}},
  ['Taml'] = {{0x0B80, 0x0BFF}},
  ['Telu'] = {{0x0C00, 0x0C7F}},
  ['Tfng'] = {{0x2D30, 0x2D7F}},
  ['Thai'] = {{0x0E00, 0x0E7F}},
  ['Tibt'] = {{0x0F00, 0x0FFF}},
  ['Vaii'] = {{0xA500, 0xA63F}},
  ['Yiii'] = {{0xA000, 0xA48F}, {0xA490, 0xA4CF}}
}

Babel.script_blocks.Cyrs = Babel.script_blocks.Cyrl
Babel.script_blocks.Hant = Babel.script_blocks.Hans
Babel.script_blocks.Kana = Babel.script_blocks.Jpan

\end{lstlisting}



\clearpage

\ifinswedish
\printglossary[style=mylong, type=\acronymtype, title={Lista över akronymer
och förkortningar}]
\else
%\printglossary[style=mylong, type=\acronymtype, title={List of Acronyms and abbreviations}]
%\printnoidxglossaries% Replace \printglossary[type=\acronymtype, ...] with:
\printnoidxglossary[type=\acronymtype, style=mylong, title={List of Acronyms and Abbreviations}]
\fi
\clearpage

% --- print the bibliography ---
% Print the bibliography (and make it appear in the table of contents)

% Specify the title of the references page
\renewcommand{\bibname}{References}

%\typeout{Biblatex current language is \currentlang}
\printbibliography[heading=bibintoc, title={References}]


\end{document}

